\begin{textsample}{POS Dim 3 – vem\_ed\_25 – Score 7.00 – vem\_ed\_25/t134.txt}  \label{ex:f3_pos_022}
4º Congresso Red LatAM COIL: \textbf{diversidade}, equidade e \textbf{inclusão} Workshop apresentado por Osvaldo Succi Junior (CPS / Brasil), José Luis Jiménez (Universidad Católica Andrés Bello, Venezuela) e Mirjam Hauck (The Open University, Reino Unido) Entre 4 e 6 de setembro de 2024, ocorreu o 4º Congresso da Red LatAM COIL, com o tema “\textbf{diversidade}, equidade e \textbf{inclusão}”.
Um dos objetivos da rede é expandir os benefícios da abordagem COIL (Collaborative Online International Learning) como estratégia de internacionalização do currículo na educação superior, envolvendo países da América Latina e outras partes do mundo.
Sessões paralelas, \textbf{painéis} de discussão e workshops compuseram o evento on-line, que contou com as convidadas especiais Hope Windle (SUNY COIL Center) e Rosi León (DePaul University).
Na abertura, após a palavra da presidenta da Red LatAM COIL, Brenda García Portillo, apresentaram-se os \textbf{membros} do \textbf{conselho} executivo - entre eles, Ana Carolina Freschi, da equipe dos PCIs / Cesu.
O tema da conferência de abertura proferida por Hope Windle, diretora do SUNY COIL Center, foi “COIL and diversity, equity and inclusion: beyond icebergs and stereotypes to community building”.
Hope Windle ressaltou que \textbf{diversidade} começa pelo esforço dos professores em pronunciar corretamente o nome dos alunos de várias origens.
Celebrou também a translinguagem (spanglish, portuñol) e defendeu sessões de quebra-gelo (icebreakers) em que se desconstroem estereótipos.
Apresentou o acrônimo TTLC (time, trust, leadership, communication) como caminho para a \textbf{inclusão}, despertando empatia e \textbf{consideração} pela \textbf{diversidade} de pensamento.
No segundo dia da conferência, Osvaldo Succi Junior, coordenador dos PCIs / Cesu, apresentou, com José Luis Jiménez (Universidad Católica Andrés Bello, Venezuela) e Mirjam Hauck (The Open University, Reino Unido) o workshop “Humanizing STEM through COIL and Critical Virtual Exchange: Empowering Educators to Foster Interdisciplinary and Culturally Inclusive Learning”. continuação A proposta de humanização do STEM (Science, Technology, Engineering, Mathematics) por meio dos projetos COIL passa pelo alinhamento com os Objetivos de Desenvolvimento Sustentável (ODS) da ONU.
Três estratégias são possíveis: colaboração entre cursos STEM-STEM – exemplos: química, ajudando comunidades sul-africanas a produzir sabões artesanais para melhorar a renda familiar; matemática, analisando dados da migração venezuelana para compreender os desafios dos migrantes; desenvolver projetos que envolvam disciplinas de humanidades e ciências “duras” para abordar desafios sociais – exemplo: ciência da computação e jornalismo, desenvolvendo jogos que despertam consciência sobre paz, justiça e direitos humanos; envolver ONGs dos dois países participantes do projeto colaborativo e propor soluções para as comunidades – exemplo: projeto entre universidades do Brasil e da Venezuela que colaboraram com ONGs para lidar com a remoção forçada de comunidades indígenas devido à extração mineral e ao \textbf{desmatamento} na Bacia Amazônica.
Osvaldo Succi Junior sugere, nesse contexto de humanização dos projetos, uma nova proposta de PCI / COIL.
Em vez de desenhar o projeto com base em disciplinas, partir de um ou mais ODS para a solução de problemas locais ou globais.
Na parte prática do workshop, os participantes esboçaram um projeto COIL contemplando o ODS 3 (Saúde e Bem-Estar) e o 10 (Redução das Desigualdades).
Tarefa principal: desenvolver campanhas de comunicação para desconstruir mitos relacionados às vacinas no México e nas Bahamas, em colaboração com ONGs como Médicos sem Fronteiras e Cruz Vermelha.

% matched lemmas: conselho, consideração, desmatamento, diversidade, inclusão, membro, painel
\end{textsample}
